\section{The \texttt{Selector} McStas Component}
velocity selector (helical lamella type) such as \textless{}b\textgreater{}V\_selector\textless{}/b\textgreater{} component

\subsection*{Identification}
\begin{itemize}
  \item \textbf{Site:} 
  \item \textbf{Author:} Peter Link, \textless{}a href="mailto:Andreas.Ostermann@frm2.tum.de"\textgreater{}Andreas Ostermann\textless{}/a\textgreater{}
  \item \textbf{Origin:} Uni. Gottingen (Germany)
  \item \textbf{Date:} MARCH 1999
\end{itemize}

\subsection*{Description}
\begin{lstlisting}
Velocity selector consisting of rotating Soller-like blades
defining a helically twisted passage.
Geometry is defined by two identical apertures at 12 o'clock position,
The origin is at the ENTRANCE of the selector.

Example: Selector(xmin=-0.015,  xmax=0.015,  ymin=-0.025, ymax=0.025, length=0.25,
nslit=72,d=0.0004, radius=0.12, alpha=48.298, nu=500)
These are values for the D11@ILL Dornier 'Dolores' Velocity Selector (NVS 023)

%VALIDATION
Jun 2005: extensive external test, one minor problem
Jan 2006: problem solved (for McStas-1.9.1)
Validated by: K. Lieutenant

%BUGS
for transmission calculation, each neutron is supposed to be in the guide centre
\end{lstlisting}

\subsection*{Input parameters}
Parameters in \textbf{boldface} are required; the others are optional.

\begin{longtable}{p{0.22\textwidth}p{0.12\textwidth}p{0.46\textwidth}p{0.14\textwidth}}
\toprule
\textbf{Name} & \textbf{Unit} & \textbf{Description} & \textbf{Default} \\
\midrule
\endhead
xmin & m & Lower x bound of entry aperture & -0.015 \\
xmax & m & Upper x bound of entry aperture & 0.015 \\
ymin & m & Lower y bound of entry aperture & -0.025 \\
ymax & m & Upper y bound of entry aperture & 0.025 \\
length & m & rotor length & 0.25 \\
xwidth & m & Width of entry. Overrides xmin,xmax. & 0 \\
yheight & m & Height of entry. Overrides ymin,ymax. & 0 \\
nslit & 1 & number of absorbing subdivinding spokes/lamella & 72 \\
d & m & width of spokes at beam-center & 0.0004 \\
radius & m & radius of beam-center & 0.12 \\
alpha & deg & angle of torsion & 48.298 \\
nu & Hz & frequency of rotation, which is ideally 3956*alpha*DEG2RAD/2/PI/lambda/length & 500 \\
\bottomrule
\end{longtable}

\subsection*{Links}
\begin{itemize}
  \item \href{run:/home/willend/willend-McCode/mcstas-comps/optics/Selector.comp}{Source code} for \texttt{Selector.comp}.
  \item See also Additional notes \textless{}a href="http://mcstas.risoe.dk/pipermail/neutron-mc/1999q1/000134.html"\textgreater{}March 1999\textless{}/a\textgreater{} and \textless{}a href="http://mcstas.risoe.dk/pipermail/neutron-mc/1999q2/000136.html"\textgreater{}Jan 2000\textless{}/a\textgreater{}.
\end{itemize}
\IfFileExists{Selector_static.tex}{\input{Selector_static.tex}}{}